\chapter[Measure and probability vocabulary]{The measure-theoretic vocabulary behind probability}
\label{ch:measure-vocabulary}
A probability model is not merely a formula for averages. It states
\emph{which outcomes exist}, \emph{which events can be assigned
probabilities}, and \emph{which calculations are allowed}.
This chapter provides the missing hypotheses for Results 21--24,
40--46, and 71--76.

\begin{definition}[Sigma-algebra and measure]
For a set $\Omega$, a $\sigma$-algebra $\mathcal F$ is a family of
subsets containing $\Omega$, closed under complements and countable
unions. A measure $\mu:\mathcal F\to[0,\infty]$ satisfies
$\mu(\varnothing)=0$ and countable additivity over disjoint events.
A probability measure has $\mu(\Omega)=1$. A measure is
\emph{$\sigma$-finite} if $\Omega$ is a countable union of
measurable sets of finite measure.
\end{definition}
\begin{definition}[Random variable and expectation]
A real-valued random variable $X$ is a measurable map
$(\Omega,\mathcal F)\to(\mathbb R,\mathcal B(\mathbb R))$:
for every Borel set $B$, its preimage $X^{-1}(B)$ is an event.
If $\int|X|\,d\Prob<\infty$, define $\E X=\int X\,d\Prob$.
The space $L^2$ consists of equivalence classes of variables with
$\E|X|^2<\infty$; equality almost surely is identified.
Its inner product is $\langle X,Y\rangle=\E[X\overline Y]$.
A \emph{Hilbert space} is a complete inner-product space; $L^2$
is one. Completeness is what makes orthogonal projection possible.
\end{definition}
\begin{proposition}[Conditional expectation is a projection]
For $X\in L^2$ and a sub-$\sigma$-algebra $\mathcal G\subseteq
\mathcal F$, the conditional expectation $\E[X\mid\mathcal G]$
is characterized by being $\mathcal G$-measurable and satisfying
$\E[(X-\E[X\mid\mathcal G])Z]=0$ for all bounded
$\mathcal G$-measurable $Z$. It is the orthogonal projection of
$X$ onto the closed subspace $L^2(\mathcal G)$.
\end{proposition}
\begin{proof}
The defining conditional-expectation identity gives orthogonality
for bounded $Z$. Truncation followed by $L^2$ approximation
extends it to $Z\in L^2(\mathcal G)$.
The projection theorem then characterizes the minimizer of
$\E[(X-Z)^2]$ among such $Z$.
\end{proof}
This explains why a conditional firing-rate prediction is an
optimal mean-square prediction \emph{relative to the variables
conditioned upon}, not relative to every unobserved brain state.

\begin{definition}[Stationarity, ergodicity, and predictability]
A process $X_t$ is \emph{wide-sense stationary} if its mean is
constant and $\Cov(X_t,X_s)$ depends only on $t-s$.
It is \emph{strictly stationary} if all finite-dimensional
joint laws are time-shift invariant. An invariant process is
\emph{ergodic} when invariant events have probability zero or one.
A predictable intensity $\lambda(t\mid\mathcal H_{t-})$ is
measurable relative to information available \emph{just before}
$t$; this forbids predicting a spike using its own occurrence.
The time-integrated intensity is called its \emph{compensator}
when the Doob--Meyer conditions are satisfied.
For a regular parametric model $p_\theta$, the \emph{score} is
$U_\theta=\nabla_\theta\log p_\theta(X)$ and
\emph{Fisher information} is
$\mathcal I(\theta)=\E_\theta[U_\theta U_\theta^\top]$,
under assumptions justifying its finiteness and calculus.
\end{definition}
\begin{proposition}[Law of total variance revisited]
For $X\in L^2$,
\[
\Var(X)=\E[\Var(X\mid\mathcal G)]
 +\Var(\E[X\mid\mathcal G]).
\]
\end{proposition}
\begin{proof}
Write $X-\E X=(X-M)+(M-\E X)$ where
$M=\E[X\mid\mathcal G]$. The cross term vanishes by
orthogonality. Taking squared expectations yields the identity.
\end{proof}
\begin{example}[LFP recording versus ensemble]
If a stationary and ergodic model describes a long LFP trace,
its time average can estimate the ensemble mean under integrability.
Stationarity alone does not justify equating those means, nor does
increasing the sample rate create independent trials.
\end{example}
% --- expanded edition (October 2026) ---
\section{Modes of convergence and their counterexamples}
For random variables on one probability space, four modes of convergence
are in common use:
\begin{itemize}[leftmargin=2em]
\item \emph{almost sure}: $\Prob(X_n\to X)=1$;
\item \emph{in $L^p$}: $\E|X_n-X|^p\to0$;
\item \emph{in probability}: $\Prob(|X_n-X|>\varepsilon)\to0$ for every $\varepsilon>0$;
\item \emph{in distribution}: $\Prob(X_n\le x)\to\Prob(X\le x)$ at continuity points of the limit.
\end{itemize}
Almost sure convergence and $L^p$ convergence each imply convergence in
probability, which implies convergence in distribution. No other
implication holds in general.
\begin{proposition}
$L^p$ convergence implies convergence in probability.
\end{proposition}
\begin{proof}
Markov's inequality applied to $|X_n-X|^p$ gives
$\Prob(|X_n-X|>\varepsilon)\le\E|X_n-X|^p/\varepsilon^p\to0$.
\end{proof}
\begin{example}[Convergence without convergence of means]
On $\Omega=(0,1)$ with Lebesgue measure let $X_n=n\ind_{(0,1/n)}$.
Then $X_n\to0$ almost surely, but $\E X_n=1$ for all $n$. The
spike-count analogue is a rare, large burst: it vanishes in probability
but keeps a fixed contribution to the mean.
\end{example}
\begin{example}[Convergence in probability without almost sure convergence]
Enumerate the dyadic intervals $[0,1],[0,\tfrac12],[\tfrac12,1],[0,\tfrac14],\ldots$
and let $X_n$ be the indicator of the $n$th one. The lengths tend to
zero, so $X_n\to0$ in probability, but every $\omega$ lies in
infinitely many of the intervals, so $X_n(\omega)$ does not converge.
\end{example}
Consistency of an estimator is usually stated in probability; strong
consistency is almost sure. A central limit theorem is a statement in
distribution, and says nothing about any particular realization.

\section{Product measures and independence}
Independence of $X$ and $Y$ means that their joint law is the product
of the marginals: $\Prob_{(X,Y)}=\Prob_X\otimes\Prob_Y$. For $\sigma$-finite
measures the product measure is unique on the product $\sigma$-algebra,
and the Tonelli and Fubini theorems (Chapter~\ref{ch:classical-sampling})
license computing expectations of functions of independent variables
by iterated integration. Independence of a whole family requires the
product property for every finite subfamily.

\section{Filtrations, stopping times, and martingales}
A \emph{filtration} $(\mathcal F_t)$ is an increasing family of
$\sigma$-algebras representing information available at time $t$.
A process $M_t$ adapted to it, with $\E|M_t|<\infty$, is a
\emph{martingale} if $\E[M_t\mid\mathcal F_s]=M_s$ for $s\le t$.
\begin{example}[Counting-process martingale]
For a homogeneous Poisson process of rate $\lambda$ with its natural
filtration, $M_t=N(t)-\lambda t$ is a martingale: for $s<t$,
$\E[N(t)-N(s)\mid\mathcal F_s]=\lambda(t-s)$ by independent
increments. More generally $N(t)-\int_0^t\lambda(u\mid\mathcal H_{u-})\,du$
is a (local) martingale for the conditional intensity, which is
Result 74.
\end{example}
A \emph{stopping time} $\tau$ is a random time such that
$\{\tau\le t\}\in\mathcal F_t$: whether it has happened is known from
the past. The first spike after stimulus onset is a stopping time;
``the spike that precedes the largest ripple of the session'' is not.
\begin{theorem}[Optional stopping, bounded case]
If $M$ is a martingale and $\tau$ a bounded stopping time, then
$\E M_\tau=\E M_0$.
\end{theorem}
\begin{proof}[Proof in discrete time]
If $\tau\le K$, write $M_\tau=M_0+\sum_{k=1}^K(M_k-M_{k-1})\ind_{\{\tau\ge k\}}$.
The event $\{\tau\ge k\}=\{\tau\le k-1\}^c$ belongs to $\mathcal F_{k-1}$,
so each term has expectation
$\E[\ind_{\{\tau\ge k\}}\E(M_k-M_{k-1}\mid\mathcal F_{k-1})]=0$.
\end{proof}
Applied to $N(t)-\lambda t$, optional stopping gives
$\E N(\tau)=\lambda\E\tau$ for bounded stopping times: the expected
number of spikes before a history-dependent stopping rule equals the
rate times the expected duration, as long as the rule does not look
into the future. Analyses that end an epoch ``at the last spike'' violate
this and are biased.

\section*{Exercises for this chapter}
\begin{exercise}
Show that the intersection of two $\sigma$-algebras is a $\sigma$-algebra
but their union need not be.
\end{exercise}
\begin{exercise}
Show that convergence in distribution to a constant implies convergence
in probability to that constant.
\end{exercise}
\begin{exercise}
For a Poisson process, show that $(N(t)-\lambda t)^2-\lambda t$ is a
martingale.
\end{exercise}
\begin{exercise}
Decide whether ``the first time the LFP envelope exceeds five MADs'' and
``the time of the session's largest envelope value'' are stopping times.
\end{exercise}
